\documentclass[aspectratio=169]{beamer}
\usetheme{simple}

\input{preamble/preamble.tex}
\input{preamble/preamble_math.tex}
% \input{preamble/preamble_acronyms.tex}

\title{Lecture 4: Generalized Linear Models}
\subtitle{STATS305B: Applied  Statistics II}
\author{Scott Linderman}
\date{\today}


\begin{document}


\maketitle
    
\begin{frame}{Recap}

Last time...
\begin{itemize}
    \item Definition and Examples of Exponential Family Distributions
    \item The Log Normalizer
    \item Maximum Likelihood Estimation
    \item Mean Parameterization
    \item KL Divergence and Deviance Residuals
\end{itemize}
    
\end{frame}
    
\begin{frame}{Outline}

    Today...
    \begin{itemize}
        \item Generalized Linear Models (GLMs)
        \item Canonical form
        \item Demo: Poisson GLM
        \item Non-canonical forms
        \item Model checking and comparison
    \end{itemize}
    
\end{frame}
    
\begin{frame}[t,allowframebreaks]{Generalized Linear Models}\label{generalized-linear-models}

Logistic regression was a special case of a more general class of models
called \textbf{generalized linear models} (GLMs). 

In a GLM, the conditional distribution \(p(Y \mid X=x)\) is modeled as an exponential
family distribution whose mean parameter is a function of \(X\). 

For example, if \(Y \in \naturals\), we could model it with a Poisson GLM;
if \(Y \in \{1,\ldots,K\}\), we could model it as a categorical GLM. 

Many of the nice properties of logistic regression carry over!
\end{frame}

\begin{frame}[t,allowframebreaks]{Model}\label{model}

To construct a generalized linear model with exponential family
observations, we set \begin{align*}
    \E[y_i \mid \mbx_i] &= f(\mbbeta^\top \mbx_i).
\end{align*}

From above, this implies, \begin{align*}
    \nabla A(\eta_i) &= f(\mbbeta^\top \mbx_i) \\
    \Rightarrow \eta_i &= [\nabla A]^{-1} \big( f(\mbbeta^\top \mbx_i) \big),
\end{align*} when \(\nabla A(\cdot)\) is invertible. (In this case, the
exponential family is said to be \textbf{minimal}).

The \textbf{canonical mean function} is \(f(\cdot) = \nabla A(\cdot)\)
so that \(\eta_i = \mbbeta^\top \mbx_i\).

The (canonical) \textbf{link function} is the inverse of the (canonical)
mean function.
\end{frame}

\begin{frame}[t,allowframebreaks]{Logistic regression
revisited}\label{logistic-regression-revisited}

Consider the Bernoulli distribution once more. The gradient of the log
normalizer is, \begin{align*}
    \nabla A(\eta) &= \nabla \log (1 + e^\eta) 
    = \frac{e^\eta}{1+ e^\eta}
\end{align*} This is the logistic function!

Thus, logistic regression is a Bernoulli GLM with the canonical mean
function.
\end{frame}

\begin{frame}[t,allowframebreaks]{Example: Poisson GLM}\label{example-poisson-glm}

Recall that the Poisson distribution can be written in exponential
family form as, \begin{align*}
\mathrm{Po}(y; \lambda) 
= \frac{1}{y!} \lambda^y e^{-\lambda} 
= \frac{1}{y!} \exp \left\{ y \log \lambda - \lambda  \right\} 
&= h(y) \exp \left\{ \langle t(y), \eta \rangle - A(\eta) \right\} 
\end{align*} where 
\begin{itemize}
    \item \(h(y) = 1/y!\)
    \item \(t(y) = y\)
    \item \(\eta = \log \lambda\) 
    \item \(A(\eta) = e^\eta\)
\end{itemize}

The canonical mean function is
\(f(x^\top \beta) = \nabla A(x^\top \beta) = e^{x^\top \beta}\). 

The canonical mean function is the exponential.
\end{frame}

\begin{frame}[t,allowframebreaks]{Log likelihood: Canonical case}\label{canonical-case}

Canonical mean functions lead to nice math. Consider the log joint
probability, \begin{align*}
    \cL(\mbbeta) 
    &= \sum_{i=1}^n \langle t(y_i), \eta_i \rangle - A(\eta_i)  + c \\
    &= \sum_{i=1}^n \langle t(y_i), \mbbeta^\top \mbx_i \rangle - A(\mbbeta^\top \mbx_i) + c,
\end{align*} where we have assumed a canonical mean function so
\(\eta_i = \mbbeta^\top \mbx_i\).
\end{frame}

\begin{frame}[t,allowframebreaks]{Gradient of the log
likelihood: Canonical case}\label{gradient-of-the-log-likelihood}

The gradient is, \begin{align*}
    \nabla \cL(\mbbeta) 
    &= \sum_{i=1}^n \langle t(y_i), \mbx_i \rangle - \langle \nabla A(\mbbeta^\top \mbx_i), \, \mbx_i \rangle\\
    &= \sum_{i=1}^n \langle t(y_i) - \E[t(Y); \eta_i], \, \mbx_i \rangle
\end{align*} where \(\eta_i = \mbx_i^\top \mbbeta\) for a GLM with
canonical link.

In many cases, \(t(y_i) = y_i \in \reals\) so \begin{align*}
    \nabla \cL(\mbbeta) 
    &= \sum_{i=1}^n (y_i - \hat{y}_i) \mbx_i.
\end{align*}
\end{frame}

\begin{frame}[t,allowframebreaks]{Hessian of the log
likelihood: Canonical case}\label{hessian-of-the-log-likelihood}

When \(t(y_i) = y_i\), the Hessian is \begin{align*}
    \nabla^2_{\mbbeta} \cL(\mbbeta) 
    &= - \sum_{i=1}^n \nabla^2 A(\mbbeta^\top \mbx_i) \, \mbx_i \mbx_i^\top \\
    &= -\sum_{i=1}^n \Var[Y; \eta_i] \, \mbx_i \mbx_i^\top
\end{align*}
\end{frame}

\begin{frame}[t,allowframebreaks]{Newton updates}\label{newton-updates}

Now recall the undamped Newton's method updates, written here in terms
of the change in weights, \begin{align*}
    \Delta \mbbeta &= - [\nabla^2 \cL(\mbbeta)]^{-1} \nabla \cL(\mbbeta) \\
    &= \left[\sum_{i=1}^n \Var[t(Y); \eta_i] \, \mbx_i \mbx_i^\top \right]^{-1} \left[\sum_{i=1}^n (y_i - \hat{y}_i) \mbx_i \right]
\end{align*}

Letting \(w_i = \Var[t(Y); \eta_i]\), \begin{align*}
    \Delta \mbbeta &=
    \left[\sum_{i=1}^n w_i \, \mbx_i \mbx_i^\top \right]^{-1} \left[ \sum_{i=1}^n (y_i - \hat{y}_i) \mbx_i \right] \\
    % &= (\mbX^\top \mbW \mbX)^{-1} [\mbX^\top \mbW \mbW^{-1} (\mby - \hat{\mby})] \\
    &= (\mbX^\top \mbW \mbX)^{-1} [\mbX^\top \mbW \hat{\mbz}]
\end{align*} where \(\mbW = \diag([w_1, \ldots, w_i])\) and
\(\hat{\mbz} = \mbW^{-1} (\mby - \hat{\mby})\).
\end{frame}

\begin{frame}[t,allowframebreaks]{Iteratively reweighted least
squares}\label{iteratively-reweighted-least-squares}

This is \textbf{iteratively reweighted least squares (IRLS)} with
weights \(w_i\) and working responses \begin{align*}
\hat{z}_i = \frac{y_i - \hat{y}_i}{w_i},
\end{align*} both of which are functions of the current weights
\(\mbbeta\).

As in logistic regression, the working responses can be seen as linear
approximations to the observed responses mapped back through the link
function.
\end{frame}

\begin{frame}[allowframebreaks]{Demo: Neural Spike Train
Analysis}\label{demo-neural-spike-train-analysis}

\begin{center}
    \textit{Go to Colab notebook}
\end{center}
\end{frame}

\begin{frame}[t,allowframebreaks]{Non-canonical case}\label{non-canonical-case}

When we choose an arbitrary mean function, the expressions are a bit
more complex. Let's focus on the case where \(t(y_i) = y_i\) for scalar
\(y_i\), but allow for arbitrary mean function \(f\). \begin{align*}
    \cL(\mbbeta) 
    &= \sum_{i=1}^n \langle y_i, \eta_i \rangle - A(\eta_i)  + c,
\end{align*} but now
\(\eta_i = [\nabla A]^{-1} f(\mbbeta^\top \mbx_i)\).

The gradient is, \begin{align*}
    \nabla \cL(\mbbeta) 
    &= \sum_{i=1}^n \left(y_i - \hat{y}_i\right) \frac{\partial \eta_i}{\partial \mbbeta}.
\end{align*}

Applying the inverse function theorem, as above, yields, \begin{align*}
\frac{\partial \eta_i}{\partial \mbbeta} 
&= \mathrm{Var}[Y]^{-1} \mbx_i = \mbx_i / w_i,
\end{align*} and \begin{align*}
    \nabla \cL(\mbbeta) 
    &= \sum_{i=1}^n \left(\frac{y_i - \hat{y}_i}{w_i} \right) f'(\mbbeta^\top \mbx_i) \mbx_i.
\end{align*}

With automatic differentiation, it's not a big deal to fit GLMs with non-canonical mean functions using gradient descent or Newton's method. 
The math just isn't as clean.

\end{frame}

\begin{frame}[t,allowframebreaks]{Deviance and Goodness of
Fit}\label{deviance-and-goodness-of-fit}

We can perform maximum likelihood estimation via Newton's method, but do
the resulting parameters \(\hat{\mbbeta}_{\mathsf{MLE}}\) provide a good
fit to the data? 

One way to answer this question is by comparing the
fitted model to two reference points: a \textbf{saturated} model and a
\textbf{baseline} model.

For binomial GLMs (including logistic regression) and Poisson GLMs, the
saturated model conditions on the mean equaling the observed response.

For example, in a Poisson GLM the saturated model's log probability is,
\begin{align*}
\log p_{\mathsf{sat}}(\mby) 
&= \sum_{i=1}^n \log \mathrm{Po}(y_i; y_i) \\
&= \sum_{i=1}^n -\log y_i! + y_i \log y_i - y_i.
\end{align*}

The likelihood ratio statistic in this case is \begin{align*}
-2 \log \frac{p(\mby \mid \mbX; \hat{\mbbeta})}{p_{\mathsf{sat}}(\mby)} 
&= 2 \sum_{i=1}^n \log \mathrm{Po}(y_i; y_i) - \log \mathrm{Po}(y_i; \hat{\mu}_i)\\
&= 2 \sum_{i=1}^n y_i \log \frac{y_i}{\hat{\mu}_i} + \hat{\mu}_i - y_i \\
&= \sum_{i=1}^n r_{\mathsf{D}}(y_i, \hat{\mu}_i)^2
\end{align*} 
where \(\hat{\mu}_i = f(\mbx_i^\top \hat{\mbbeta})\) is the
predicted mean. 

We recognize the likelihood ratio statistic as the sum
of squared deviance residuals! (See last lecture.)

Moreover, this statistic is just the deviance (twice the KL divergence)
between the two models, \begin{align*}
2 \KL{p_{\mathsf{sat}}(\mbY)}{p(\mbY \mid \mbX=\mbx; \hat{\mbbeta})}
&= 2 \E_{p_{\mathsf{sat}}} \left[\log \frac{p_{\mathsf{sat}}(\mby)}{p(\mby \mid \mbX; \hat{\mbbeta})} \right] \\
&= 2 \sum_{i=1}^n \KL{\mathrm{Po}(y_i)}{\mathrm{Po}(\hat{\mu}_i)} \\
&\triangleq D(\mby, \hat{\mbmu}).
\end{align*} 

The same idea holds for GLMs with other exponential family
distributions. 

Larger deviance implies a poorer fit.
\end{frame}

\begin{frame}[t,allowframebreaks]{Deviance and Goodness of
    Fit: Baseline Model}\label{deviance-and-goodness-of-fit}
    
Larger deviance relative to what?

The baseline model is typically a GLM with only an intercept term, in
which case the MLE is
\(\mu_i \equiv \frac{1}{n} \sum_{i=1}^n y_i = \bar{y}\). 

For that
baseline model, the deviance is, \begin{align*}
D(\mby, \bar{y} \mbone)
&= 2 \sum_{i=1}^n y_i \log \frac{y_i}{\bar{y}} + \bar{y} - y_i \\
&= 2 \sum_{i=1}^n y_i \log \frac{y_i}{\bar{y}} 
\\
&= \sum_{i=1}^n r_{\mathsf{D}}(y_i, \bar{y})^2.
\end{align*}

\end{frame}


\begin{frame}[t,allowframebreaks]{Fraction of Deviance
Explained}\label{fraction-of-deviance-explained}

As with linear models where we use the fraction of variance explaiend
(\(R^2\)), in GLMs we can consider the fraction of \emph{deviance}
explained. 

Note that the deviance is positive for any model that is not
saturated. 

The \textbf{fraction of deviance explained} is
\begin{align*}
    1 - \frac{D(\mby; \hat{\mbmu})}{D(\mby; \bar{y} \mbone)}
\end{align*} 

Unless your model is worse than guessing the mean, the fraction of deviance
explained is between 0 and 1.
\end{frame}

\begin{frame}[t,allowframebreaks]{Model Comparison}\label{model-comparison}

The difference in deviance between two models with predicted means
\(\hat{\mbmu}_0\) and \(\hat{\mbmu}_1\), the difference in deviance,
\begin{align*}
D(\mby; \hat{\mbmu}_0) - D(\mby; \hat{\mbmu}_1)
\end{align*} has an approximately chi-squared null distribution. We can
use this fact to sequentially add or subtract features from a model
depending on whether the change in deviance is significant or not.
\end{frame}

\begin{frame}[t,allowframebreaks]{Model Checking}\label{model-checking}

Just like in standard linear models, we should inspect the residuals in
generalized linear models for evidence of model misspecification. For
example, we can plot the residuals as a function of \(\hat{\mu}_i\) and
they should be approximately normal at all levels of \(\hat{\mu}_i\).
\end{frame}

\begin{frame}[t,allowframebreaks]{Cross-Validation}\label{cross-validation}

Finally, another common approach to model selection and comparison is to
use cross-validation. The idea is to approximate out-of-sample
predictive accuracy by randomly splitting the training data.

\textbf{Leave-one-out cross validation (LOOCV)} withholds one datapoint
out at a time, estimates parameters using the other \(n-1\), and then
evaluates predictive likelihood on the held out datapoint.

This approximates, \begin{align}
    \E_{p^\star(x, y)}[\log p(y \mid x, \{x_i, y_i\}_{i=1}^n)] &\approx
    \frac{1}{n} \sum_{i=1}^n \log p(y_i \mid x_i, \{x_j, y_j\}_{j \neq i}).
\end{align} where \(p^\star(x, y)\) is the true data generating
distribution

For small \(n\) (or when using a small number of folds), a
bias-correction can be used.

% Cross-validated predictive log likelihood estimates are similar to the
% \emph{jackknife} resampling method in the sense that it is estimating an
% expectation wrt the unknown data-generating distribution \(p^\star\) by
% resampling the given data.

\end{frame}

\begin{frame}[t,allowframebreaks]{Conclusion}\label{conclusion}

Generalized linear models allow us to build flexible regression models that respect the domain of the response variable. 

Logistic regression is a special case of a Bernoulli GLM with the canonical link function. 

For categorical data, we could use a categorical distribution with the
softmax link function, and for count data, we could use a Poisson GLM
with exponential, softplus, or other link functions. 

Leveraging the deviance and deviance residuals for exponential family distribiutions,
we can derive analogs of familiar terms from linear modeling, like the
fraction of variance explained and the residual analysis.

\end{frame}
    
    
\end{document}
